\subsection{Модуль lab1}

\subsubsection{GraphMetrics: эксцентриситеты}

\textbf{Вход:} \texttt{Graph const\&} (конструктор).

\textbf{Выход:} \texttt{MetricsResult}~--- \texttt{unordered\_map<int,double>}
эксцентриситетов, \texttt{vector<int>} центра и диаметральных вершин,
\texttt{double} радиус и диаметр.

\textbf{Описание:} Граф ациклический и невзвешенный с точки зрения метрик:
эксцентриситет считается по числу рёбер. Поэтому  используется BFS~--- приватный метод \texttt{bfsSteps},
который возвращает расстояния в шагах от источника. Сложность одного
обхода $O(V + E)$, итоговая~--- $O(V(V+E))$.

После обхода \texttt{compute} вычисляет эксцентриситет каждой вершины
как максимум расстояний по всем остальным. $+\infty$ означает, что
вершина недостижима.

\begin{lstlisting}[style=cppstyle, caption={BFS и вычисление метрик (GraphMetrics.cc)}]
unordered_map<int, double> GraphMetrics::bfsSteps(int src) const
{
    constexpr double INF = numeric_limits<double>::infinity();
    std::unordered_map<int, double> dist;
    for (int v : m_graph_.vertexIds()) dist[v] = INF;
    dist[src] = 0.0;

    std::queue<int> q;
    q.push(src);
    while (!q.empty()) {
        int u = q.front(); q.pop();
        for (auto const& [v, w] : m_graph_.neighbors(u)) {
            if (dist[v] == INF) {
                dist[v] = dist[u] + 1.0;
                q.push(v);
            }
        }
    }
    return dist;
}

MetricsResult GraphMetrics::compute() const {
    constexpr double INF = std::numeric_limits<double>::infinity();
    auto ids = m_graph_.vertexIds();
    MetricsResult result;
    double radius = INF, diameter = 0.0;

    for (int v : ids) {
        auto dist = bfsSteps(v);
        double ecc = 0.0;
        for (int u : ids) {
            if (u == v) continue;
            double d = dist.at(u);
            if (d == INF) { ecc = INF; break; }
            ecc = std::max(ecc, d);
        }
        result.eccentricities[v] = ecc;
        if (ecc != INF) {
            radius   = std::min(radius, ecc);
            diameter = std::max(diameter, ecc);
        }
    }
    result.radius   = radius == INF ? 0.0 : radius;
    result.diameter = diameter;
    for (int v : ids) {
        double ecc = result.eccentricities[v];
        if (ecc == result.radius)   result.center.push_back(v);
        if (ecc == result.diameter) result.diametralVerts.push_back(v);
    }
    std::sort(result.center.begin(), result.center.end());
    std::sort(result.diametralVerts.begin(), result.diametralVerts.end());
    return result;
}
\end{lstlisting}

\subsubsection{ShimbellMethod: матричная степень}

\textbf{Вход:} \texttt{Graph const\&} (конструктор), \texttt{int pathLength}.

\textbf{Выход:} \texttt{ShimbellResult} --- два поля имеют тип вида
\texttt{DistanceMatrix}~$=$~\\\texttt{vector<vector<optional<double>>>}:
\texttt{min\_distances} и \texttt{max\_distances}, и \texttt{int path\_length}.

\textbf{Описание:} Начальная матрица строится через вспомогательный метод\\
\texttt{CollectionUtils::makeSquareMatrix} с геттером-лямбдой:
диагональ~--- нули, ребро~--- его вес из \texttt{getEdgeWeight},
иначе~--- \texttt{nullopt}. Матрица итеративно умножается $k-1$ раз.
Умножение вынесено в приватные методы \texttt{multiplyMin} и
\texttt{multiplyMax}, которые делегируют в
\texttt{CollectionUtils::multiplyOptionalMatrix} с соответствующим
компаратором. Использование \texttt{nullopt}
исключает ложные пути через несуществующие рёбра. Краевой случай $k=0$ возвращает единичную матрицу.

\begin{lstlisting}[style=cppstyle, caption={Матрица смежности и метод Шимбелла (ShimbellMethod.cc)}]
DistanceMatrix ShimbellMethod::createAdjacencyMatrix() const {
    return CollectionUtils::makeSquareMatrix<optional<double>>(
        m_vertex_ids_,
        [&](int u, int v) -> std::optional<double> {
            if (u == v) return 0.0;
            auto w = m_graph_.getEdgeWeight(u, v);
            return w.has_value() ? w : std::nullopt;
        }
    );
}

ShimbellResult const& ShimbellMethod::compute(int pathLength) {
    if (pathLength == 0) {
        DistanceMatrix identity(m_size_,
            vector<optional<double>>(m_size_, nullopt)
        );
        for (int i = 0; i < m_size_; ++i) identity[i][i] = 1;
        result_ = ShimbellResult{identity, identity, 0};
        return result_;
    }

    DistanceMatrix current_min = createAdjacencyMatrix();
    DistanceMatrix current_max = createAdjacencyMatrix();

    DistanceMatrix base_min    = current_min;
    DistanceMatrix base_max    = current_max;

    for (int step = 2; step <= pathLength; ++step) {
        current_min = multiplyMin(current_min, base_min);
        current_max = multiplyMax(current_max, base_max);
    }
    result_ = ShimbellResult{
        .min_distances = current_min,
        .max_distances = current_max,
        .path_length   = pathLength
    };
    return result_;
}

DistanceMatrix ShimbellMethod::multiplyMin(DistanceMatrix const& a, DistanceMatrix const& b) const {
    return CollectionUtils::multiplyOptionalMatrix(
    a, b, std::less<>()
    );
}

DistanceMatrix ShimbellMethod::multiplyMax(DistanceMatrix const& a, DistanceMatrix const& b) const {
    return CollectionUtils::multiplyOptionalMatrix(
    a, b, std::greater<>()
    );
}
\end{lstlisting}

\subsubsection{PathCounter: подсчёт маршрутов}

\textbf{Вход:} \texttt{Graph const\&} (конструктор), \texttt{int from},
\texttt{int to}.

\textbf{Выход:} \texttt{PathMatrix}~$=$~\texttt{vector<vector<int>>} --- все пути от \texttt{from} до \texttt{to}.

\textbf{Описание:} Публичный метод \texttt{getAllPaths} сначала проверяет
наличие обеих вершин через \texttt{hasVertex} и возвращает пустой результат
при отсутствии, не запуская обход. Затем вызывается рекурсивный
\texttt{backtrackAllPaths}~--- DFS с маской посещённых вершин
\texttt{unordered\_map<int,bool>}. При достижении цели текущий путь
копируется в результат; при выходе из рекурсии вершина снимается с
отметки (\emph{backtracking}). Сложность: $O(n!)$ в худшем случае.

\begin{lstlisting}[style=cppstyle, caption={Backtracking-поиск маршрутов (PathCounter.cc)}]
PathMatrix PathCounter::getAllPaths(int from, int to) {
    if (!m_graph_.hasVertex(from) || !m_graph_.hasVertex(to))
        return {};
    PathMatrix all_paths;
    std::vector<int> current_path;
    std::unordered_map<int, bool> visited;
    backtrackAllPaths(from, to, visited,
                      current_path, all_paths
    );
    return all_paths;
}

void PathCounter::backtrackAllPaths(
    int current, int target,
    std::unordered_map<int, bool>& visited,
    std::vector<int>& currentPath, PathMatrix& allPaths)
{
    currentPath.push_back(current);
    visited[current] = true;
    if (current == target) {
        allPaths.push_back(currentPath);
    } else {
        for (auto const& [neighborId, _] : m_graph_.neighbors(current))
            if (!visited[neighborId])
                backtrackAllPaths(neighborId, target, visited,
                                  currentPath, allPaths);
    }
    currentPath.pop_back();
    visited[current] = false;
}
\end{lstlisting}
